% Databento-only AAO reduced-form article; prospective analysis, NOT run results.
% This paper uses no Barchart option marks or historical CL data.
% No fabricated address, email, experimental count or inference.
\documentclass[a4paper,fleqn]{cas-dc}
\usepackage[authoryear,longnamesfirst]{natbib}
\usepackage{amsmath,amssymb,booktabs}
\usepackage{graphicx}

\begin{document}
\let\WriteBookmarks\relax
\shorttitle{Volatility uncertainty in average-price options}
\shortauthors{Espino Montelongo}
\title[mode=title]{Historical and Implied Volatility in WTI Average Price Options:
A Bayesian Assessment Using Exchange Settlement Data}
\author[1]{Heriberto Espino Montelongo}
\affiliation[1]{organization={Universidad de las Am\'ericas Puebla},
 city={San Andr\'es Cholula},state={Puebla},country={Mexico}}

\begin{abstract}
We develop an empirical design for assessing whether integrating historical
volatility uncertainty changes the valuation of WTI average-price options
relative to more consequential differences among volatility inputs.
The evidence is drawn solely from exchange contract definitions and
settlement statistics provided through Databento. We estimate the
risk-neutral expected terminal average and a discount factor using
call--put parity at reserved strikes, and price other contracts under
a reduced-form lognormal approximation to the terminal arithmetic
average. Volatility estimates use only prior-date option-implied
forward-average changes or preceding implied-volatility observations.
We compare posterior integration, posterior-mean plug-in pricing,
two rolling estimates, EWMA and lagged option-implied volatility on
identical option-date observations. We distinguish actual-flagged
and theoretical-flagged settlements and use date-cluster inference.
The structural first-nearby WTI futures model cannot be identified
from the purchased AAO definitions and settlement statistics alone;
the proposed analysis is therefore a cross-sectional, reduced-form
exercise, not a fully independent ex-ante futures-based pricing test.
\textit{Empirical results will be inserted only after the complete
reproducible Databento run and its validation.}
\end{abstract}

\begin{keywords}
Asian options \sep WTI crude oil \sep Bayesian inference \sep
historical volatility \sep implied volatility \sep exchange settlements
\end{keywords}
\maketitle

\section{Introduction}
\label{sec:db_intro}

Parameter uncertainty is distinct from volatility-input selection in
option pricing. Under nonlinear valuation, averaging option values
over a posterior distribution of volatility need not equal evaluating
the same pricing function at posterior-mean volatility. The difference
depends on price curvature and posterior dispersion. Conversely, even
an accurate posterior under the physical distribution need not identify
risk-neutral volatility; pricing-model assumptions and volatility risk
premia remain relevant \citep{guidolin2003,romboutsstentoft2014}.

We examine these differences for WTI average-price options using
one source of observed exchange settlement statistics.
Unlike the earlier Barchart-based study, this study does not use
Barchart option histories or independent LO option marks.
The study separates the question of \emph{posterior integration}
from the empirical comparison of historical and previously observed
same-family implied-volatility inputs.

CME's WTI Average Price Option is a European-style option on the
arithmetic average of daily first-nearby CL futures settlements
during its contract month. This contract specification implies a
multi-fixing payoff, which existing Asian-option methods model
explicitly \citep{kemnavorst1990,curran1994,kahale2020}.
The purchased Databento data here comprise AAO definitions
and option settlement statistics, not the CL term structure or the
underlying daily fixing history. We therefore estimate the
\emph{risk-neutral expected average} from AAO put--call parity and
make a transparent lognormal approximation to the terminal average.
We do not present that approximation as exact structural Asian pricing.

\section{Data and eligibility}
\label{sec:db_data}

The proposed data window comprises January 2020 through September
2026, subject to availability and final integrity checks after
the archive's reacquisition. We identify contracts by publisher,
native symbol, expiry, strike and option right, rather than by
temporary numeric instrument ID. Settlement messages are reconciled
across source-file months in publication timestamp order, including
preliminary reports, later finalizations, and correction/delete
actions. A contract--reference-date mark is eligible only if its
latest reconstructed action is active, final, non-intraday and
has a finite, nonnegative price.

A settlement marked \emph{actual} by the exchange is not necessarily
a transaction execution, and a \emph{theoretical} settlement may be
model-derived. We retain this distinction in all empirical summaries
and do not count millions of correlated option-date marks as
independent statistical replications. Model results are organized by
distinct CME settlement reference date.

We restrict the principal pricing sample to valuation dates
\emph{before the first calendar day of the option's averaging month}.
This deliberately avoids claiming to know already-realized CL daily
fixings, which are absent from the purchased AAO data. Prices and
strikes are required to satisfy separately documented unit
sanity checks. The data-eligibility funnel, final number of
dates, final contract-date count and method common support
are \textit{pending actual execution}.

\section{Pricing design}
\label{sec:db_method}

\subsection{Expected average and discount from separate parity strikes}

For a European arithmetic-average call and put sharing strike $K$,
contract expiry and reference date, the parity identity is
\begin{equation}
 C_t(K)-P_t(K)
 =D(t,T)\{\mathbb{E}_t^Q[A_T]-K\}.
 \label{eq:db_parity}
\end{equation}
Write $F_{A,t}=\mathbb{E}_t^Q[A_T]$ and
fit, across a predesignated set of matched call--put
\emph{anchor strikes}, the cross-sectional relation
\begin{equation}
 C_t(K)-P_t(K)=a_t-b_t K,\qquad
 D(t,T)=b_t,\quad F_{A,t}=a_t/b_t.
 \label{eq:db_reg}
\end{equation}
A group requires at least four available paired strikes and
sensible estimated discount and forward values. Groups that
fail parity-dispersion, strike-unit or discount checks are excluded.
\textbf{Every anchor strike is removed from both sides of the
evaluated target cross-section.} This avoids fitting an option's
own price mechanically through parity. Nevertheless, information
from \emph{other} current-date options determines the same-date
forward average; the exercise is a conditional cross-sectional
valuation, not a wholly prior-date prediction.

\subsection{Lognormal terminal-average approximation}

Because the first-nearby CL daily fixing path is not observed in
this dataset, we use the reduced-form approximation
\begin{equation}
 A_T \approx F_{A,t}
  \exp\!\left(-\tfrac12\sigma_A^2\tau+
  \sigma_A\sqrt{\tau}\,Z\right),\quad Z\sim N(0,1),
\end{equation}
where $\tau$ is time to option expiry and $\sigma_A$ is a
\emph{proxy average-forward volatility}, not the original
first-nearby CL spot/futures diffusion coefficient.
Its discounted call price is
\begin{equation}
 C_t^B(\sigma_A)=D(t,T)
 \bigl[F_{A,t}\Phi(d_1)-K\Phi(d_2)\bigr],\quad
 d_{1,2}=\frac{\log(F_{A,t}/K)}{\sigma_A\sqrt{\tau}}
 \pm \tfrac12\sigma_A\sqrt{\tau}.
\end{equation}
Put prices follow parity. This compact model provides a controlled
comparison of volatility estimators under a common pricing rule;
it does not reconstruct the exact arithmetic-average payoff law.

\subsection{Historical proxy volatility and Bayesian pricing}

For each expiry, define the \emph{option-implied forward-average}
return using parity estimates on successive eligible earlier
reference dates:
\begin{equation}
 r_u=\log F_{A,u}-\log F_{A,u-1}.
\end{equation}
We use only available past returns, omitting large gaps, and
annualize by a 252-observation convention. These are NOT
observations of first-nearby crude-oil futures returns and
cannot by themselves estimate the latter's physical volatility.
For $y_u=\sqrt{252}\,r_u$, let
$y_u\mid m,v\sim N(m,v)$ and
$v\sim\mathrm{InvGamma}(\alpha_0,\beta_0)$, with
$(\alpha_0,\beta_0)=(3,0.16)$ and the usual integrated
flat-prior treatment of the unknown mean. The resulting shape
and scale are
\begin{equation}
 \alpha_n=\alpha_0+(n-1)/2,\quad
 \beta_n=\beta_0+\tfrac12\sum_u(y_u-\bar y)^2.
\end{equation}
Draw $\sigma_A=\sqrt v$ from the posterior. We compare
posterior-integrated and posterior-mean plug-in values,
\begin{equation}
 C_{PI}=\mathbb E[C_t^B(\sigma_A)\mid\mathcal D_{<t}],
 \qquad C_{PM}=C_t^B(\mathbb E[\sigma_A\mid\mathcal D_{<t}]).
\end{equation}
Their second-order difference is approximately
$\tfrac12 C_t^{B\prime\prime}(\bar\sigma_A)
\operatorname{Var}(\sigma_A\mid\mathcal D_{<t})$,
subject to the smoothness and concentration of the posterior.

\subsection{Benchmark volatility estimates and timing}

We compare six specifications on identical contract-date rows:
posterior integration (PI), posterior-mean plug-in (PM),
21- and 63-observation rolling volatility, exponentially weighted
volatility with 63-observation half-life, and prior-date
AAO-implied volatility. The last is obtained by inverting the
same reduced-form pricing function on earlier non-anchor
option marks, selecting near-the-money contracts.
Every volatility input is determined strictly before $t$;
only the disjoint-strike parity estimate of $F_{A,t}$ and
$D(t,T)$ uses the current-date AAO surface.

\section{Empirical evaluation}
\label{sec:db_results}

We evaluate all six methods on their common target contracts
and reference dates. For method $m$, pricing errors are
$e_{i,t}^{(m)}=\widehat C_{i,t}^{(m)}-C_{i,t}^{\mathrm{settlement}}$,
and report bias, mean absolute error and root mean squared error.
The principal comparisons are paired at the date level;
valuation dates are resampled in the cluster bootstrap so
contract-level dependence is not mistaken for independent evidence.

We separately report actual- and theoretical-flagged settlement
marks, number of eligible dates, and the fraction of attempted
expiry-date groups lost at the parity, volatility-history,
lagged-IV or pricing stages. A comparison favoring a prior
option-implied specification would establish reduced-form
same-family information content, not an independently identified
variance risk premium or a claim about actual execution prices.

\subsection{Results to be inserted after execution}

\textit{The Databento-only experiment has not yet been run and validated.
No benchmark rankings, confidence intervals, statistical significance
or out-of-sample error figures are asserted here. The evaluation
runner will produce method\_metrics.csv, daily\_metrics.csv,
cluster\_intervals.csv and an experiment manifest.}

\IfFileExists{figures/fig06_databento_benchmark_mae.pdf}{
\begin{figure*}[t]
\centering
\includegraphics[width=.9\textwidth]{figures/fig06_databento_benchmark_mae.pdf}
\caption{Common-support reduced-form valuation MAE by settlement-mark type;
generated only after running the empirical analysis.}
\end{figure*}
}{}
\IfFileExists{figures/fig07_databento_date_cluster_intervals.pdf}{
\begin{figure*}[t]
\centering
\includegraphics[width=.82\textwidth]{figures/fig07_databento_date_cluster_intervals.pdf}
\caption{Paired date-cluster improvements relative to posterior-mean
plug-in pricing; confidence intervals are based on date resampling.}
\end{figure*}
}{}
\IfFileExists{figures/fig08_databento_mae_time_series.pdf}{
\begin{figure*}[t]
\centering
\includegraphics[width=.9\textwidth]{figures/fig08_databento_mae_time_series.pdf}
\caption{Reduced-form method errors through time; points summarize
date-level errors rather than independent contract-level samples.}
\end{figure*}
}{}

\section{Limitations and robustness}
\label{sec:db_limits}
This design has three nontrivial limitations. First, parity requires
current-day prices for other AAO contracts; selection of non-overlapping
anchor and target strikes prevents direct price reuse but does not
create independent ex-ante volatility forecasts. Second, the
historical series consists of parity-implied average-forward levels,
not actual returns on the constituent first-nearby WTI futures.
Third, theoretical CME settlements may be generated from pricing
procedures related to the evaluated benchmarks.

Sensitivity analysis will re-estimate the comparison on the
actual-flagged subsample and under alternative parity-anchor
sizes, tighter parity-residual cutoffs and posterior-prior
choices. The numerical assumptions and complete common-support
funnel must be archived before journal submission.

\section{Conclusion}
\label{sec:db_conclusion}
The study is designed to distinguish the incremental numerical
effect of integrating historical proxy-volatility uncertainty from
the choice of volatility source using one exchange-supplied AAO
dataset. Its scope is necessarily narrower than structural Asian
pricing with observed futures-fixing histories. Conclusions on
model ranking and economic magnitude are deliberately deferred
until a fully audited Databento-only run has been completed.

\section*{Data availability}
The study relies on Databento CME AAO contract definitions and
settlement statistics. The privately maintained research archive,
source hashes and reproducibility manifests record the exact
sample vintage. Public redistribution is subject to the vendor's
terms; no additional Barchart option data are used.

\section*{Code availability}
The code and an explicit execution recipe are maintained in
\texttt{github.com/heri-espino/Bayesian-Uncertainty-in-WTI-APOs}.
A final version will identify a frozen source commit and empirical run.

\section*{Declaration of competing interests}
The author declares no competing interests.

\bibliographystyle{cas-model2-names}
\bibliography{references}
\end{document}
